\section{About the Committee}

\subsection{History of the Committee}

UNODC was established in 1997, formed through the merger of the United Nations International Drug Control Programme (UNDCP), created in 1991, and the Crime Prevention and Criminal Justice Division of the United Nations Office at Vienna, established in 1992. It was originally named the Office for Drug Control and Crime Prevention and adopted its current name in 2002 to reflect its broadened scope.\footnote{United Nations Office on Drugs and Crime, "About UNODC" (UNODC, no date) \url{https://www.unodc.org/unodc/en/about-unodc/index.html}}

Since its founding, UNODC's mandate has expanded considerably beyond drug control, shaped by a series of landmark developments:

\begin{guidetable}[Y{0.27}Y{1.73}]{2}
  \textbf{Year} & \textbf{Event} \\
  \midrule
  1997 & Establishment of the Office for Drug Control and Crime Prevention through the merger of UNDCP and the Crime Prevention and Criminal Justice Division. \\
  2000 & Adoption of the United Nations Convention against Transnational Organized Crime (UNTOC) in Palermo, together with the supplementing Protocol to Prevent, Suppress and Punish Trafficking in Persons, Especially Women and Children. This is the primary legal instrument underpinning this Committee's mandate on trafficking.\footnote{United Nations General Assembly, "United Nations Convention against Transnational Organized Crime" (res 55/25, 15 November 2000). The text of the Convention and Protocol is available via UNODC: \url{https://www.unodc.org/unodc/en/organized-crime/intro/UNTOC.html}} \\
  2002 & The Office adopted its current name, the United Nations Office on Drugs and Crime, reflecting its widened focus on organised crime alongside drug control. \\
  2003 & Adoption of the United Nations Convention against Corruption (UNCAC), for which UNODC also serves as guardian and secretariat. \\
  2016-201 & The Security Council formally linked human trafficking to armed conflict and terrorism for the first time through Resolution 2331 (2016), reinforced by Resolution 2388 (2017), both calling on UNODC and member states to strengthen victim identification, prosecution and cross-border cooperation in conflict-affected regions.\footnote{United Nations Security Council, Resolution 2331 (2016), S/RES/2331 (20 December 2016) \url{https://main.un.org/securitycouncil/en/s/res/2331-\%282016\%29}} \footnote{United Nations Security Council, Resolution 2388 (2017), S/RES/2388 (21 November 2017) \url{https://main.un.org/securitycouncil/en/content/sres2388-2017}} \\
  2010-20 & Leadership of the Office has reflected its growing international profile: Yuri Fedotov served as Executive Director from 2010 to 2019, followed by Ghada Fathi Waly of Egypt from 2020 to 2026. On 11 May 2026, Monica Juma of Kenya assumed office as Executive Director of UNODC and Director-General of the United Nations Office at Vienna, becoming the most senior Kenyan official in United Nations history.\footnote{United Nations, "Secretary-General Appoints Monica Juma of Kenya Executive Director, United Nations Office on Drugs and Crime" (UN Press Release, 2026). The specific press release should be available at \url{https://press.un.org}} \\
\end{guidetable}

\subsection{Mandate of the UNODC}

The United Nations Office on Drugs and Crime (UNODC) is the entity within the UN Secretariat mandated to assist member states in their struggle against illicit drugs, transnational organised crime, corruption and terrorism in all their forms and manifestations. Its mission, in line with the current UNODC Strategy, is to contribute to global peace and security, human rights and development by making the world safer from drugs, crime and terrorism, working for and with member states to promote justice, the rule of law and resilient societies.\footnote{United Nations Office on Drugs and Crime, UNODC Strategy 2021–2025 (Vienna: UNODC, 2021) \url{https://www.unodc.org/unodc/en/strategy/index.html}}

UNODC pursues this mission through three interlocking pillars of work:

\begin{itemize}
  \item \textbf{Research and analysis}: producing evidence-based knowledge products, such as the annual Global Report on Trafficking in Persons and the World Drug Report, which inform policy-making at the national and international level.
  \item \textbf{Normative work}: acting as guardian of, and secretariat to, the treaty bodies overseeing the United Nations Convention against Transnational Organized Crime (UNTOC) and its supplementing Protocols, including the Palermo Protocol on trafficking in persons, as well as the United Nations Convention against Corruption (UNCAC) and the three international drug control conventions.
  \item \textbf{Technical assistance and capacity-building}: training law enforcement, judicial and border officials, supporting legislative reform, and strengthening victim identification and protection systems, particularly in states with limited institutional capacity.\footnote{United Nations Office on Drugs and Crime, UNODC Strategy 2021–2025 (Vienna: UNODC, 2021) \url{https://www.unodc.org/unodc/en/strategy/index.html}}
\end{itemize}

As a specialised, non-enforcement body of the UN Secretariat, UNODC's mandate is exercised through cooperation, technical support and normative guidance rather than binding enforcement measures. Delegates should therefore frame resolutions around treaty implementation, institutional capacity-building, information-sharing and victim protection, rather than measures, such as sanctions or peacekeeping deployment, that fall outside the Committee's competence.

Because trafficking in persons frequently overlaps with organised crime, terrorism financing and conflict dynamics, UNODC exercises this mandate in close coordination with other UN entities, including UNHCR, IOM and UNICEF, as well as with the Security Council when trafficking intersects with threats to international peace and security.\footnote{United Nations Office on Drugs and Crime, Global Report on Trafficking in Persons 2024 (Vienna: UNODC, 2024) \url{https://www.unodc.org/unodc/en/data-and-analysis/glotip.html}}
